\documentclass{article}
\usepackage[utf8]{inputenc}
\usepackage[T1]{fontenc}
\usepackage[french]{babel}
\usepackage{amsmath}
\usepackage{mathtools}
\usepackage{amssymb}

\begin{document}

\section{Exercice 4 : Mathématiques et LaTeX}

\subsection{Questions 1 à 6 : Bases, modes et indices}
% Q1-Q2-Q3
Mode normal : 3-2+4=5. \\
Mode math : $3-2+4=5$.

% Q4
Accent circonflexe pour les exposants : en mode normal $2^{10}=1024$ (ou affiché séparément), en mode math $2^{10} = 1024$.

% Q5-Q6
Variables indicées et produits : 
$a_1, a_2, \dots, a_n$ et $a_1 \times a_2 \times \cdots \times a_n$.

\subsection{Questions 7 à 9 : Inégalités, symboles et paquets}
% Q7-Q8
Inégalités et relations : $3 < 4$, $5 > 4$, $3 \leq 4$, $5 \geq 4$. \\
Quantificateurs et relations : $\exists$, $\forall$, $\neq$, $\equiv$, $\iff$, $\sim$, $\hookrightarrow$. \\
Fonctions et fractions : $\cos(x)$, $\sin(x)$, $\tan(x)$, $\lim_{x \to 0}$, $\frac{a}{b}$, $\infty$.

% Q9
Symboles avancés (paquets amssymb / mathtools) : 
$3 \nleq 4$, $3 \ngeq 4$, $A \implies B$.

\subsection{Questions 10 à 14 : Équations en ligne à part et étiquettes}
% Q10-Q11
\begin{equation*}
    e = mc^2
\end{equation*}

% Q12-Q13-Q14
\begin{equation}
    \label{eq:einstein}
    E = mc^2
\end{equation}
Référence avec \verb|\eqref| : Comme vu dans l'équation \eqref{eq:einstein}.

\subsection{Questions 15 à 18 : Texte, quantificateurs et limites complexes}
% Q15
\begin{equation}
    \forall k < n, \text{ $k$ est un diviseur de } n!
\end{equation}

% Q16 : Cours-TD sur les quantificateurs (Exemple d'énoncé et solution)
\paragraph{Énoncé (Q16) :} Démontrer que pour tout réel $x$, si $x > 2$ alors $x^2 > 4$.
\begin{equation*}
    \forall x \in \mathbb{R}, \quad (x > 2 \implies x^2 > 4)
\end{equation*}

% Q17
\begin{equation}
    \lim_{\theta \to \left(\frac{\pi}{2}\right)^-} \tan(\theta) = \infty
\end{equation}

\end{document}